% Generated by code/make_appendix_g.py -- do not edit by hand.
\begin{table}[htbp]
\centering
\caption{Every model trained for the second round, with the compute it consumed. GPU seconds are wall-clock on one RTX~4060 Laptop (8\,GB).}
\label{tab:gpu-ledger}
\footnotesize
\setlength{\tabcolsep}{4pt}
\begin{tabular}{llrrrrl}
\toprule
Run & Model & Params & Ep. & Sd. & GPU\,s & Test \\
\midrule
A (baseline) & TCN, 4 blocks, 96 ch & 243,873 & 18 & 1 & 16.5 & AUC 0.8220 / AP 0.0676 \\
B (dual window) & two 96-ch TCN branches & 487,617 & 18 & 1 & 57.0 & AUC 0.8323 / AP 0.1065 \\
C (dual + aux heads) & two branches, two heads & 512,387 & 18 & 1 & 54.3 & AUC 0.8345 / AP 0.1076 \\
P2 quantile GRU & GRU 80$\to$96, 2 layers & 206,753 & 40 & 3 & 42.1 & 3-D RMSE 1.080 cm at 0.2 s \\
P3 reservation learner & 2$\times$128 MLP, REINFORCE & 23,054 & 200 & 1 & 402.7 & S2 +5.0 / S3 -6.9 pt, fail \\
\midrule
Total & & & & & 572.6 & \\
\bottomrule
\end{tabular}
\end{table}

\begin{table}[htbp]
\centering
\caption{P1-b detector, three arms on the official environment split. The pass mark was fixed before the second and third arms were trained: test AUC $\ge 0.85$ or test AP $\ge 0.09$, with the per-environment floor as a veto.}
\label{tab:p1b-arms}
\footnotesize
\begin{tabular}{lrrrrr}
\toprule
Arm & Val AUC & Val AP & Test AUC & Test AP & Worst env.\ AUC \\
\midrule
A baseline & 0.8853 & 0.1121 & 0.8220 & 0.0676 & 0.7646 \\
B dual window & 0.8663 & 0.1096 & 0.8323 & 0.1065 & 0.7865 \\
C dual + multi-task & 0.8798 & 0.1131 & 0.8345 & 0.1076 & 0.7884 \\
\bottomrule
\end{tabular}
\end{table}

\begin{table}[htbp]
\centering
\caption{The same detector on three seeds. The pre-registered bar had three parts: average precision up at least fifteen per cent, every seed above $0.09$, and no more than two points lost on the worst environment. The first two hold; the third misses by four hundredths of a point out of a two-point allowance, which is why the upgrade is reported as rejected rather than accepted.}
\label{tab:p1b-seeds}
\footnotesize
\begin{tabular}{rrrrrr}
\toprule
Seed & A test AP & C test AP & A worst env. & C worst env. & $\Delta$ worst \\
\midrule
20261120 & 0.0676 & 0.1076 & 0.7646 & 0.7884 & 0.0238 \\
20261122 & 0.0692 & 0.0918 & 0.7189 & 0.7299 & 0.0110 \\
20261123 & 0.0786 & 0.1554 & 0.7345 & 0.7154 & -0.0191 \\
\midrule
Mean & 0.0718 & 0.1182 & 0.7393 & 0.7446 & 0.005260 \\
\bottomrule
\end{tabular}
\end{table}

\begin{table}[htbp]
\centering
\caption{P2 quantile predictor against the ridge replica at four mask lengths. Errors are three-dimensional RMSE in centimetres, mean over three seeds (standard deviation in the second column). The last column is the paired difference over the seven held-out records; a positive number means the quantile model is closer.}
\label{tab:p2quant}
\footnotesize
\begin{tabular}{rrrrrrr}
\toprule
Mask $L$ & Quantile RMSE & s.d. & Ridge RMSE & Quantile $q_{90}$ & Coverage & Paired $\Delta$ [95\,\% CI] \\
\midrule
6 & 1.080 & 0.003 & 2.006 & 1.598 & 0.9122 & 0.802 [0.561, 1.044] \\
18 & 4.028 & 0.011 & 4.526 & 6.825 & 0.9114 & 0.363 [0.126, 0.600] \\
30 & 6.140 & 0.130 & 6.314 & 9.968 & 0.8967 & 0.214 [0.119, 0.310] \\
\bottomrule
\end{tabular}
\end{table}

\begin{table}[htbp]
\centering
\caption{The vectorised scheduler engine and what the extra speed bought. Left: bit-exact agreement with the reference implementation and the measured speed-up. Right: the paid-review arm over the full grid, judged against a pre-declared one-point bar.}
\label{tab:p3engine}
\footnotesize
\begin{tabular}{lrr}
\toprule
Quantity & Reference & Vectorised \\
\midrule
One candidate, 30 seeds (s) & 6.438 & 0.220 \\
Thirty candidates (s) & 26.38 & 0.126 \\
960-candidate search unit (s) & 844.0 & 13.4 \\
Speed-up, worst case & \multicolumn{2}{r}{29.2$\times$} \\
Speed-up, 30-candidate batch & \multicolumn{2}{r}{209.2$\times$} \\
Speed-up, search unit & \multicolumn{2}{r}{63.0$\times$} \\
\midrule
Bit-exact agreement, max $|\Delta|$ & \multicolumn{2}{r}{0.0e+00} \\
\bottomrule
\end{tabular}
\end{table}

% numerals-ok: begin
\noindent\textbf{What the detector upgrade bought.} The baseline arm reproduces the published detector to $10^{-6}$ on both splits, so the comparison is against the same model the body reports. Widening the input window doubles the parameters and lifts average precision from 0.0676 to 0.1065; adding the two auxiliary heads lifts it once more, to 0.1076 on the seed the architecture search used. Three further seeds put the same comparison on a firmer footing and also break it: average precision rises by 64.7 per cent against the same-seed baseline (75.0 per cent against the pre-registered single-seed figure) and every seed clears $0.09$, but the worst-environment score falls by 0.53 hundredths of a point more than the two-point allowance, so the pre-registered verdict is a rejection. The honest reading is that the ranking signal improved and the calibration did not: at a $0.5$ cut the detector loses 9.5 per cent of the good frames to catch 44.6 per cent of the bad ones, which is why Part~II never takes the flag as an input.

\noindent\textbf{What the quantile predictor bought.} At the shortest mask the quantile model halves the error of the ridge replica (1.080 against 2.006~cm). At the thirty-frame mask it does not: the pre-declared bar of 5.80~cm is missed by 0.340~cm. What survives is the decision, not the regression: because the objective $J(T)$ prices review interval against error, the better short-horizon model moves the optimum from 0.30 to 0.25~s and lowers the objective by 19.3 per cent (4.5212 to 3.6486). A model that is worse on the long horizon is still the better model to deploy.

\noindent\textbf{What a paid second look bought.} Nothing, on this task. Across 90 cells (fifteen case--mapping combinations, six reservation levels, thirty seeds each) the review arm fires in only 35 of 90 cells and never clears the one-point bar: the best cell is 0.278 points and the worst is -0.556. The information is real --- the second prefix moves the predicted finish time by about 0.70~s on average --- but the value-of-information gate fires in one trial in twenty, and the time a review buys comes out of the same budget it is meant to protect, so the net effect on the on-time rate sits inside the noise. That is the sharpest form of the result in Section~\ref{sec:limits}: the observation exists, is paid for, and is correctly scheduled, and the decision is still unchanged. The engine that makes this measurable is the vectorised simulator, which replays the reference bit-for-bit and runs the whole grid in 31.6~s.

\noindent\textbf{What the video bought.} The full re-extraction decodes every released video: 1047 windows, two camera streams, 1{,}047 aligned training rows' worth of frames. Adding the resulting 184 summary features to the 94 pose features lowers test MAE from 3.1942 to 2.4352~s, which looks like a gain until the permutation control lands at 2.4258~s. The gain is the shrinkage of a wider design matrix, not information: permuting the video columns changes nothing, so the video block is not used by the model and is not used by the paper's decisions.

\noindent\textbf{What a learned reservation bought.} Nothing that survives its own bar. A REINFORCE policy with 23,054 parameters and thirteen reservation levels, trained on fifteen seeds for 402.7 GPU-seconds over three runs, lifts S2 by 5.00 points and drops S3 by 6.88 on the five held-out seeds: the mean of -0.94 points is below the pre-declared bar of $+2.00$ with both interval lower bounds above zero, so the verdict is a failure. The failure is structural rather than a test-slice accident: the same pairing already sits below zero on the ten validation seeds, the weighted rate moves in lumps of $1/W$ with weights between one and three, two or three of the five seeds are unchanged, and the paired 95 per cent intervals are 15.3 and 14.9 points wide --- wider than the effect in both cases. The grid corner $q=0.30$ --- the best fixed reserve on validation --- gives S2 $+6.7$ and S3 $+0.6$ points on the same seeds, so no reserve-only policy in this family reaches the bar on both cases. What the two arms do show is an interaction: the risk selector learns a state-dependent reserve (S2 at level 1, S3 at level 0) while the slack selector collapses to a constant (level 3 in 88--100 per cent of decisions). A reserve that only shifts cost and feasibility, never the completion time, is a liability the data cannot price --- Proposition~B, recovered as a training result.
% numerals-ok: end

